\section{Solution Method}
\label{sec:solution}

\subsection{Reduced System}
\label{sec:solution_reduced}

The Taylor rule~\eqref{eq:nk-taylor} determines the nominal interest
rate algebraically as a contemporaneous function of inflation, output,
and the policy disturbance. Substituting this rule into the dynamic IS
curve~\eqref{eq:nk-is} yields a system with two genuinely
forward-looking jumpers ($y_{t}$ and $\pi_{t}$) and three exogenous
AR(1) state variables ($\varepsilon_{t}^{a},\varepsilon_{t}^{b},
\varepsilon_{t}^{m}$). Stacking the variables as
$\mathbf{y}_{t}=(\varepsilon_{t}^{a},\varepsilon_{t}^{b},
\varepsilon_{t}^{m},y_{t},\pi_{t})'$, the reduced system is written
in matrix form as
\begin{equation}
  A\,\mathbb{E}_{t}\bigl[\mathbf{y}_{t+1}\bigr]=B\,\mathbf{y}_{t},
  \label{eq:matrix-form}
\end{equation}
with the coefficient matrices
\begin{equation*}
  A=
  \begin{pmatrix}
    1 & 0 & 0 & 0 & 0 \\
    0 & 1 & 0 & 0 & 0 \\
    0 & 0 & 1 & 0 & 0 \\
    0 & 0 & 0 & 1 & 1/\sigma \\
    0 & 0 & 0 & 0 & \beta
  \end{pmatrix},
  \qquad
  B=
  \begin{pmatrix}
    \rho_{a} & 0 & 0 & 0 & 0 \\
    0 & \rho_{b} & 0 & 0 & 0 \\
    0 & 0 & \rho_{m} & 0 & 0 \\
    0 & -1 & 1/\sigma & 1+\phi_{y}/\sigma & \phi_{\pi}/\sigma \\
    1 & 0 & 0 & -\kappa & 1
  \end{pmatrix}.
\end{equation*}
The first three rows of each matrix encode the AR(1) transitions of
the exogenous shocks; the fourth row is the IS curve with the Taylor
rule substituted in; and the fifth row is the New Keynesian Phillips
curve.

\subsection{Generalised Schur Decomposition}
\label{sec:solution_klein}

The rational-expectations solution is computed by the generalised
Schur (QZ) decomposition of \citet{Klein2000}. The pair $(A,B)$ is
factored as $A=Q\,T\,Z^{H}$ and $B=Q\,S\,Z^{H}$ with $T,S$
upper-triangular. The generalised eigenvalues
$\lambda_{i}=T_{ii}/S_{ii}$ relate to the dynamics-matrix eigenvalues
$\mu_{i}=1/\lambda_{i}$ of $A^{-1}B$. Stability requires
$|\mu_{i}|<1$. After reordering so that stable directions come first,
the policy function is recovered as $\mathbf{j}_{t}=F\,\mathbf{x}_{t}$
where $\mathbf{j}_{t}=(y_{t},\pi_{t})'$ and
$\mathbf{x}_{t}=(\varepsilon_{t}^{a},\varepsilon_{t}^{b},
\varepsilon_{t}^{m})'$.

\subsection{Stability and Numerical Verification}
\label{sec:solution_stability}

At the calibrated parameters, the QZ decomposition returns five
generalised eigenvalues with dynamics eigenvalues
$\{0.90,\,0.50,\,0.50,\,1.154+0.253\,\mathrm{i},\,1.154-0.253\,\mathrm{i}\}$.
The first three coincide with the shock persistences
$\rho_{a},\rho_{b},\rho_{m}$ and are stable; the remaining two form a
complex-conjugate pair with modulus $1.182$ outside the unit circle.
Three stable and two unstable eigenvalues exactly match the three
predetermined states and two forward-looking jumpers, so the
Blanchard--Kahn condition of \citet{BlanchardKahn1980} is satisfied
and the rational-expectations equilibrium is unique. The numerical
solver returns the state transition
\begin{equation}
  \mathbf{x}_{t+1}=\operatorname{diag}(0.90,\,0.50,\,0.50)\,\mathbf{x}_{t},
  \label{eq:state-trans}
\end{equation}
together with the jumper policy matrix
$F\in\mathbb{R}^{2\times 3}$ that maps the three states into $y_{t}$
and $\pi_{t}$.
